\documentclass[12pt]{article}

\usepackage{amsmath}
\usepackage{german}
\usepackage{graphicx}
\DeclareMathOperator\sgn{sgn}

%\graphicspath{{/home/ritz/Teaching/fcn/fig/}}

%----------------- layout ----------------------------------------

\oddsidemargin=0in
\textwidth=6in
\topmargin=-0.7in
\textheight=9.5in
\baselineskip=16pt
\font\mittel=cmr17
\font\roman=cmr12
\font\klein=cmr9
\font\mgross=cmr10
\newfont{\CAP}{cmcsc10 scaled 900}
\pagestyle{empty}
\parindent=0mm

%----------------- macros ----------------------------------------
\newcommand{\conc}[1]{[{\rm C}]_{\text{#1}}}
\newcommand{\units}[1]{\text{#1}}
\newcommand{\K}{\text{K}}
\newcommand{\Na}{\text{Na}}
\newcommand{\Cl}{\text{Cl}}
\newcommand{\Leak}{\text{leak}}
\newcommand{\water}{\text{H}_2\text{O}}
\newcommand{\tderiv}[1]{\frac{\mathrm{d}#1}{\mathrm{d}t}}
%-----------------------------------------------------------------


\begin{document}

\parbox{2cm}{
\includegraphics[width=1.8cm]{bccnlogo.pdf}
}
\parbox{11cm}{
\begin{center}
\large HUMBOLDT-UNIVERSIT\"AT \hskip 0.1 cm ZU \hskip 0.1 cm BERLIN
\vskip 0.1 true cm
\mgross BERNSTEIN CENTER FOR COMPUTATIONAL NEUROSCIENCE
\end{center}
}
\parbox{2cm}
{
\hfill
\includegraphics[width=1.8cm]{hublogo.pdf}
}

\vskip 0.6 true cm

\leftline{\CAP Humboldt-Universit\"at zu Berlin 
\hfill Phone: 030/2093-9110}
\leftline{\CAP Philippstr. 13 House 6
              \hfill Fax: 030/2093-6771}
\leftline{\CAP \hfill webpage: http://www.bccn-berlin.de/}


\vskip 0.8 true cm
\centerline{\bf Models of Neural Systems, WS 2010/11}
\centerline{\bf Computer Practical 7}
\centerline{Solutions to hand in on: December, 13th, 2010}


\vskip 0.8 true cm
%----------- EXERCISES BEGIN ---------------------%
\begin{enumerate}
\item \textbf{Potassium channel.} The gate model was first introduced by
    Hodgkin and Huxley to describe voltage and time dependence of
    ion conductances in the squid axon. Today it is still the standard model
    of the ion current flow through transmembrane channels. One of its main
    assumptions is that the probability of opening and closing of an ion
    gate is described by the first-order kinetic equation:
    \begin{equation}
        \tderiv{n}=\alpha_n(V)(1-n)-\beta_n(V)n
        \label{eq:activation}
    \end{equation}

    where $\alpha_n(V)$ and $\beta_n(V)$ are voltage-dependent transition
    rates.

    The potassium current in the Hodgkin-Huxley model is given by:
    \begin{equation}
        I_K=\bar{g}_{K} n(V,t)^4 (V-E_K)
        \label{eq:potassium}
    \end{equation}

    where $E_K=-77\,\mathrm{mV}$ is the reversal potential and
    $\bar{g}_K=36\,\mathrm{mS/cm^2}$ is the maximum conductance. Rates
    $\alpha_n(V)$ and $\beta_n(V)$ are given by:
    \begin{equation}
        \alpha_n(V)=\frac{0.01(V+55)}{1-\exp(-0.1(V+55))},\quad
        \beta_n(V)=0.125\exp(-0.0125(V+65)),
    \end{equation}
    where $V$ is expressed in mV, and $\alpha_n$ and $\beta_n$ are both
    expressed in units of ms$^{-1}$.
    \begin{enumerate}
        \item Write Python functions defining $\alpha_n(V)$ and $\beta_n(V)$.
        \item \label{ex:stedy-state} Plot the steady-state activation
            $n_\infty(V)=\alpha_n(V)/(\alpha_n(V)+\beta_n(V))$ and the
            activation time constant $\tau_n(V)=1/(\alpha_n(V)+\beta_n(V))$ in
            a voltage range of $-150\,\mathrm{mV}\leq V \leq 150\,\mathrm{mV}$.
        \item \textit{Voltage clamp}. Calculate the current responses
            $I_K$ to voltage steps under voltage-clamp conditions:
            \begin{equation*}
                V(t)=
                \begin{cases} V_c & \text{if $t\geq2$ ms,} \\
                    -65\,\mathrm{mV} &\text{otherwise.}
                \end{cases}
            \end{equation*}
            with the initial condition $n(t=0)=0.3177$. Plot $I_K$ as a
            function of time. Repeat the simulation for different values of
            clamping voltage $V_c$ from the range  $[-100\,\mathrm{mV},
            -40\,\mathrm{mV}]$. What can be learnt from this experiment? What
            is the predicted effect of the potassium current on the membrane
            potential? Explain the obtained results referring to the plots of
            the activation variables $n_\infty$ and $\tau_n$.
        \item \textit{Current-voltage relation}. Plot the I-V relation for
            the instantaneous and steady-state potassium
            current.
    \end{enumerate}


\item \textbf{Sodium ion channel.} The sodium channel according
    to Hodgkin-Huxley model is given by:
    \begin{equation}
        I_{\Na}(V,t)=\bar{g}_{\Na}m(V,t)^3h(V,t)(V-E_{\Na})
    \end{equation}
    where $E_{\Na}=50\,\mathrm{mV}$ is the reversal potential of sodium and
    $\bar{g}_{\Na}=120\,\mathrm{mS/cm^2}$ is its maximal conductance.
    The gating variables $m$, $h$ follow first-order kinetics (see
    Equation~\ref{eq:activation}) with the following rate functions:
    \begin{align}
        \alpha_m&=0.1\frac{V+40}{1-\exp(-0.1(V+40))} &
        \beta_m&=4\exp(-0.0556(V+65)) \\
        \alpha_h&=0.07\exp(-0.05(V+65)) &
        \beta_h&=\frac{1}{1+\exp(-0.1(V+35))}
        \label{eq:Na_gating}
    \end{align}

    \begin{enumerate}
    \item Plot the steady-state activation and time constant of $m$ and $h$ as
        a function of the membrane potential. Compare the results to the
        dynamics of the $n$ variable. Which variable is activated and which
        deactivated by depolarisation? Which variable is fast/slow?
    \item Note that in contrast to the model of the persistent potassium
        current the sodium channels are described by a system of
        ordinary differential equations.
        %In order to solve it, you need to
        %modify your Euler method! First, define a function which takes as
        %arguments the instantaneous values for $m$, $h$ and returns their
        %derivatives. Next, modify your Euler method so that it accepts a
        %function returning any number of variables together with a vector of
        %initial conditions and solves the  corresponding system of
        %differential equations.
        Calculate and plot the sodium current for $V=-20\text{mV}$,
        $m(t=0)=0.0529$, $h(t=0)=0.5961$.
    \item Simulate the voltage clamp experiment for the sodium current using
        voltage steps from -65 mV to various command voltages $V_c$. Plot the
        sodium current as a function of time. Repeat for several values of
        $V_c$. Compare the results with respective plots of the potassium
        current. What are the mechanisms responsible for the sodium current
        rise and decay?
    \item Plot the instantaneous and steady-state I-V curves for sodium
        channels. Compare and discuss the results.
\end{enumerate}

\item {\bf Report.} Submit the plain source code files for all exercises
         and a report containing the source code, all plots and answers to the
         questions asked.

% \item {\bf Hodgkin-Huxley model of action potential generation}
% 
% Action potentials (aka spikes) are a dominant feature of nervous system and
% play
% crucial role in the computations performed by cortex, thalamus,
% cerebellum and associated structures.  The ionic mechanisms of action
% potential generation were described on giant squid axon by Hodgkin and
% Huxley. The complete model they proposed is as following:
% \begin{equation}
%     c_m \tderiv{V(t)}=-I_{\Na}(V(t), t) - I_{\K}(V(t), t) - I_{\Leak}(V(t),
% t)
%     \label{eq:HHM}
% \end{equation}
% where $I_{\Na}(V,t)$, $I_{\K}(V,t)$ are the sodium and potassium currents
% which we already discussed.  $I_{\Leak}(V,t)$ is a  leakage current due to
% chloride and other ios and can be well described by a passive ion flow:
% \begin{equation}
% I_{\Leak}(V,t)=\bar{g}_{\Leak}(V-E_{\Leak})
% \end{equation}
% with $\bar{g}_{\Leak}=0.3\,\mathrm{mS/cm^2}$ and
% $E_{\Leak}=-54.387\,\mathrm{mV}$.
% 
% \begin{enumerate}
% 
% 
%     \item Simulate a Hodgkin-Huxley model by numerically integrating
%       equations for $V$, $n$, $m$, $h$ (4-dimensional system of ODEs!). Take
% $c_m=1\,\mu\mathrm{F/cm^2}$
%       and as initial values take: $V=-65\,\mathrm{mV}$, $m=0.0529$,
%       $h=0.5961$, and $n=0.3177$. Make sure that the integration time step is
% short  enough to obtain a stable and accurate solution (test several values).
%   \item Use an external current with density
%       $i_e=20\,\mu\mathrm{A/cm^2}$ and plot $V$, $I_{\K}$, $I_{\Na}$ as
%       functions of time. Which currents are responsible for different
%       phases of action potential generation?
%   \item Plot the firing rate of the model as a function of $i_e$ over
%       the range from 0 to 200 $\mu$A/cm$^2$. Show that the firing rate
%       jumps discontinuously from zero to a finite value when the
%       current passes through the minimum value to produce a sustained
%       firing. (\textit{Hint:} to detect spikes test when the membrane
%       potential crosses a predefined threshold).
%   \item Apply a pulse of a negative current with
%       $i_e=-5\,\mu\mathrm{A/cm^2}$ for 5 ms followed by $i_e=0$ and
%       show what happens. Can you explain the results?
% \end{enumerate}

\end{enumerate}

%----------- EXERCISES END -----------------------%

\vfill
\centerline{\CAP Contact}
\CAP

\begin{tabular}{lll}
Urs Bergmann & Phone: 2093-8924 & Email:
urs.bergmann(at)biologie.hu-berlin.de \\
Matthias Guggenmos & & Email: matthias.guggenmos(at)bccn-berlin.de \\
Richard Kempter \hfill & Phone: 2093-8925 \hfill & Email:
r.kempter(at)biologie.hu-berlin.de \\
\end{tabular}

\end{document}
 
